% sec1_1_part2.tex — Section 1.1 continued: Implications of Strict Exogeneity,
% Strict Exogeneity in Time-Series Models, Other Assumptions of the Model

\begin{example}[continuation of Example 1.1]
\label{ex:1.3}
For the simple regression model of Example~1.1, the strict exogeneity
assumption can be written as
\[
  \E(\varepsilon_i \mid YD_1, YD_2, \ldots, YD_n) = 0.
\]
Since $\vec{x}_i = (1,\, YD_i)'$, you might wish to write the strict exogeneity
assumption as
\[
  \E(\varepsilon_i \mid 1,\, YD_1, 1,\, YD_2, \ldots, 1,\, YD_n) = 0.
\]
But since a constant provides no information, the expectation conditional on
\[
  (1,\, YD_1, 1,\, YD_2, \ldots, 1,\, YD_n)
\]
is the same as the expectation conditional on
\[
  (YD_1, YD_2, \ldots, YD_n).
\]
\end{example}

\subsection*{Implications of Strict Exogeneity}

The strict exogeneity assumption has several implications.

\begin{itemize}
\item The \emph{un}conditional mean of the error term is zero, i.e.,
\begin{equation}
  \E(\varepsilon_i) = 0 \quad (i = 1, 2, \ldots, n).
  \label{eq:1.1.8}
\end{equation}
This is because, by the Law of Total Expectations from basic probability
theory,\footnote{The Law of Total Expectations states that
$\E[\E(y \mid \vec{x})] = \E(y)$.}
$\E[\E(\varepsilon_i \mid \vec{X})] = \E(\varepsilon_i)$.

\item If the cross moment $\E(xy)$ of two random variables $x$ and $y$ is zero,
then we say that $x$ is \textbf{orthogonal} to $y$ (or $y$ is orthogonal
to~$x$).  Under strict exogeneity, the regressors are orthogonal to the error
term for \emph{all} observations, i.e.,
\[
  \E(x_{jk}\,\varepsilon_i) = 0 \quad (i,\, j = 1, \ldots, n;\; k = 1, \ldots, K)
\]
or
\begin{equation}
  \E(\vec{x}_j \cdot \varepsilon_i)
  = \begin{bmatrix}
      \E(x_{j1}\,\varepsilon_i) \\
      \E(x_{j2}\,\varepsilon_i) \\
      \vdots \\
      \E(x_{jK}\,\varepsilon_i)
    \end{bmatrix}
  = \underset{(K \times 1)}{\vec{0}}
  \quad \text{(for all } i,\, j\text{)}.
  \label{eq:1.1.9}
\end{equation}

The proof is a good illustration of the use of properties of conditional
expectations and goes as follows.

\begin{proof}
Since $x_{jk}$ is an element of $\vec{X}$, strict exogeneity implies
\begin{equation}
  \E(\varepsilon_i \mid x_{jk})
  = \E[\E(\varepsilon_i \mid \vec{X}) \mid x_{jk}] = 0
  \label{eq:1.1.10}
\end{equation}
by the Law of Iterated Expectations from probability
theory.\footnote{The Law of Iterated Expectations states that
$\E[\E(y \mid \vec{x},\, \vec{z}) \mid \vec{x}] = \E(y \mid \vec{x})$.}
It follows from this that
\begin{align*}
  \E(x_{jk}\,\varepsilon_i)
    &= \E[\E(x_{jk}\,\varepsilon_i \mid x_{jk})]
       &\quad&\text{(by the Law of Total Expectations)} \\
    &= \E[x_{jk}\,\E(\varepsilon_i \mid x_{jk})]
       &&\text{(by the linearity of conditional expectations\footnotemark)} \\
    &= 0.
\end{align*}
\footnotetext{The linearity of conditional expectations states that
$\E[f(\vec{x})\,y \mid \vec{x}] = f(\vec{x})\,\E(y \mid \vec{x})$.}
\end{proof}

The point here is that strict exogeneity requires the regressors be orthogonal
not only to the error term from the same observation (i.e.,
$\E(x_{ik}\,\varepsilon_i) = 0$ for all~$k$), but also to the error term from
the other observations (i.e., $\E(x_{jk}\,\varepsilon_i) = 0$ for all~$k$ and
for $j \neq i$).

\item Because the mean of the error term is zero, the orthogonality conditions
\eqref{eq:1.1.9} are equivalent to zero-correlation conditions.  This is
because
\begin{align*}
  \Cov(\varepsilon_i,\, x_{jk})
    &= \E(x_{jk}\,\varepsilon_i) - \E(x_{jk})\,\E(\varepsilon_i)
       &\quad&\text{(by definition of covariance)} \\
    &= \E(x_{jk}\,\varepsilon_i)
       &&\text{(since $\E(\varepsilon_i) = 0$, see \eqref{eq:1.1.8})} \\
    &= 0
       &&\text{(by the orthogonality conditions \eqref{eq:1.1.9})}.
\end{align*}
In particular, for $i = j$, $\Cov(x_{ik},\, \varepsilon_i) = 0$.  Therefore,
strict exogeneity implies the requirement (familiar to those who have studied
econometrics before) that the regressors be contemporaneously uncorrelated with
the error term.
\end{itemize}

\subsection*{Strict Exogeneity in Time-Series Models}

For time-series models where $i$ is time, the implication \eqref{eq:1.1.9} of
strict exogeneity can be rephrased as: the regressors are orthogonal to the
past, current, and future error terms (or equivalently, the error term is
orthogonal to the past, current, and future regressors).  But for most
time-series models, this condition (and \emph{a fortiori} strict exogeneity) is
not satisfied, so the finite-sample theory based on strict exogeneity to be
developed in this section is rarely applicable in time-series contexts.
However, as will be shown in the next chapter, the estimator possesses good
large-sample properties without strict exogeneity.

The clearest example of a failure of strict exogeneity is a model where the
regressor includes the \textbf{lagged dependent variable}.  Consider the
simplest such model:
\begin{equation}
  y_i = \beta\, y_{i-1} + \varepsilon_i \quad (i = 1, 2, \ldots, n).
  \label{eq:1.1.11}
\end{equation}
This is called the \textbf{first-order autoregressive model} (AR(1)).  (We will
study this model more fully in Chapter~6.)  Suppose, consistent with the spirit
of the strict exogeneity assumption, that the regressor for observation~$i$,
$y_{i-1}$, is orthogonal to the error term for~$i$ so
$\E(y_{i-1}\,\varepsilon_i) = 0$.  Then
\begin{align*}
  \E(y_i\,\varepsilon_i)
    &= \E[(\beta\, y_{i-1} + \varepsilon_i)\,\varepsilon_i]
       &\quad&\text{(by \eqref{eq:1.1.11})} \\
    &= \beta\,\E(y_{i-1}\,\varepsilon_i) + \E(\varepsilon_i^2) \\
    &= \E(\varepsilon_i^2)
       &&\text{(since $\E(y_{i-1}\,\varepsilon_i) = 0$ by hypothesis)}.
\end{align*}
Therefore, unless the error term is always zero, $\E(y_i\,\varepsilon_i)$ is
not zero.  But $y_i$ is the regressor for observation $i+1$.  Thus, the
regressor is not orthogonal to the past error term, which is a violation of
strict exogeneity.

\subsection*{Other Assumptions of the Model}

The remaining assumptions comprising the classical regression model are the
following.

\begin{assumption}[no multicollinearity]
\label{ass:1.3}
The rank of the $n \times K$ data matrix, $\vec{X}$, is $K$ with probability~$1$.
\end{assumption}

To understand Assumption~\ref{ass:1.3}, recall from matrix algebra that the
rank of a matrix equals the number of linearly independent columns of the
matrix.  The assumption says that none of the $K$ columns of the data matrix
$\vec{X}$ can be expressed as a linear combination of the other columns of
$\vec{X}$.  That is, $\vec{X}$ is of \textbf{full column rank}.  Since the $K$
columns cannot be linearly independent if their dimension is less than~$K$, the
assumption implies that $n \geq K$, i.e., there must be at least as many
observations as there are regressors.  The regressors are said to be
\textbf{(perfectly) multicollinear} if the assumption is not satisfied.  It is
easy to see in specific applications when the regressors are multicollinear and
what problems arise.

\begin{example}[continuation of Example 1.2]
\label{ex:1.4}
If no individuals in the sample ever changed jobs, then
$\mathit{TENURE}_i = \mathit{EXPR}_i$ for all~$i$, in violation of the no
multicollinearity assumption.  There is evidently no way to distinguish the
tenure effect on the wage rate from the experience effect.  If we substitute
this equality into the wage equation to eliminate $\mathit{TENURE}_i$, the wage
equation becomes
\[
  \log(\mathit{WAGE}_i)
  = \beta_1 + \beta_2 S_i + (\beta_3 + \beta_4)\,\mathit{EXPR}_i + \varepsilon_i,
\]
which shows that only the sum $\beta_3 + \beta_4$, but not $\beta_3$ and
$\beta_4$ separately, can be estimated.
\end{example}

\begin{assumption}[spherical error variance]
\label{ass:1.4}
\leavevmode
\begin{enumerate}[label=\upshape(\roman*)]
\item \emph{(homoskedasticity)}
\begin{equation}
  \E(\varepsilon_i^2 \mid \vec{X}) = \sigma^2 > 0
  \quad (i = 1, 2, \ldots, n),
  \label{eq:1.1.12}
\end{equation}%
\footnote{When a symbol (which here is $\sigma^2$) is given to a moment (which
here is the second moment $\E(\varepsilon_i^2 \mid \vec{X})$), by implication
the moment is assumed to exist and is finite.  We will follow this convention
for the rest of this book.}

\item \emph{(no correlation between observations)}
\begin{equation}
  \E(\varepsilon_i \varepsilon_j \mid \vec{X}) = 0
  \quad (i,\, j = 1, 2, \ldots, n;\; i \neq j).
  \label{eq:1.1.13}
\end{equation}
\end{enumerate}
\end{assumption}

The homoskedasticity assumption \eqref{eq:1.1.12} says that the conditional
second moment, which in general is a nonlinear function of~$\vec{X}$, is a
constant.  Thanks to strict exogeneity, this condition can be stated
equivalently in more familiar terms.  Consider the conditional variance
$\Var(\varepsilon_i \mid \vec{X})$.  It equals the same constant because
\begin{align*}
  \Var(\varepsilon_i \mid \vec{X})
    &\equiv \E(\varepsilon_i^2 \mid \vec{X})
            - \E(\varepsilon_i \mid \vec{X})^2
       &\quad&\text{(by definition of conditional variance)} \\
    &= \E(\varepsilon_i^2 \mid \vec{X})
       &&\text{(since $\E(\varepsilon_i \mid \vec{X}) = 0$
               by strict exogeneity)}.
\end{align*}

Similarly, \eqref{eq:1.1.13} is equivalent to the requirement that
\[
  \Cov(\varepsilon_i,\, \varepsilon_j \mid \vec{X}) = 0
  \quad (i,\, j = 1, 2, \ldots, n;\; i \neq j).
\]
That is, in the joint distribution of $(\varepsilon_i,\, \varepsilon_j)$
conditional on~$\vec{X}$, the covariance is zero.  In the context of
time-series models, \eqref{eq:1.1.13} states that there is no \textbf{serial
correlation} in the error term.
